\documentclass[11pt,letterpaper]{article}
\usepackage{hanzo-defense}

\doctitle{Automated Verification and Validation Framework\\with Digital Twin Environments}
\docid{HZO-2026-004}
\docversion{Version 1.0 --- February 3, 2026}
\docdate{February 2026}
\docauthors{Hanzo AI Research --- Hanzo Team}
\docorgs{Hanzo Industries\\research@hanzo.ai}

\begin{document}
\makecover

\begin{abstract}
We present an integrated verification and validation (V\&V) framework combining continuous observability, behavioral analytics, ML pipeline orchestration, digital twin environments, and formal cryptographic proofs. The observability platform ingests logs, metrics, and distributed traces into a Datastore OLAP backend, with LLM-specific telemetry tracking token usage, cost, latency, and quality metrics across 100+ AI providers. Temporal workflows orchestrate LLM evaluation, trace clustering, and sentiment analysis, while Dagster DAGs manage ETL pipelines for training data preparation. A network simulation environment supports configurable topologies with 38 registered virtual machines, enabling mission rehearsal and system qualification. The framework includes 17 formal cryptographic proofs, 28+ consensus tests, and automated CI/CD with multi-platform builds. Anomaly detection, feature flag-controlled rollout, and A/B testing with statistical significance provide continuous data-driven decision support.
\end{abstract}

\paragraph{Executive Summary.}
This is the test, evaluation, and monitoring layer that tells operators whether a complex AI and cryptographic system actually works---before it is fielded and continuously once it is. Three capabilities work together: full instrumentation, so every action the system takes is recorded and can be inspected after the fact; a high-fidelity ``digital twin'' that replicates the real deployment, so new builds and stressful conditions (jamming, link loss, hostile inputs) can be rehearsed safely before they reach the field; and automated checks---including machine-checked cryptographic proofs---that run on every change and block anything that regresses. For a program office, this is the qualification and continuous-assurance evidence a system needs to earn and keep an authorization to operate, with specialized tracking for AI workloads (cost, latency, output quality) that conventional IT monitoring does not capture. The capability is built to run where the system runs, including disconnected at the edge with store-and-forward when connectivity returns.

\tableofcontents
\newpage

% ============================================================================
\section{Introduction}
% ============================================================================

Modern naval systems integrate AI/ML, post-quantum cryptography, and distributed computing at unprecedented scale. Traditional V\&V approaches---manual test plans, periodic audits, and waterfall certification---cannot keep pace with continuous deployment and evolving threats.

This paper presents an automated V\&V framework built on three principles:
\begin{enumerate}
\item \textbf{Continuous observability}: Every system interaction is instrumented, traced, and queryable.
\item \textbf{Automated evaluation}: ML models, cryptographic implementations, and consensus protocols are continuously verified against formal specifications.
\item \textbf{Digital twin fidelity}: System qualification occurs in high-fidelity simulated environments before live deployment.
\end{enumerate}

\subsection{Strategic Context}

A verification and validation framework is only as credible as the systems it has actually carried into production under real scrutiny. Hanzo (American-owned, Techstars~'17, founded 2014) is an applied AI lab and venture studio behind 100+ venture-funded startups and multiple unicorns; its founder is a recognized expert in \emph{both} AI/ML and cryptography, and the company builds both the \emph{Zen} frontier model family and the post-quantum cryptographic stack it runs on. The same agentic engineering stack described in these papers delivered a securities trading platform---a full ATS/BD/TA stack consolidated from 200+ microservices into 3 compiled binaries---through FINRA/SEC approval and SOC2 audit into production, in record time with a small team. That is the operative proof point for test and evaluation: regulated, externally audited software shipped at scale, the same discipline a defense accreditation demands. It also anchors a broader position. With American frontier AI consolidated into a closed two-vendor duopoly and the open frontier ecosystem now largely offshore, Hanzo is the American open-source alternative that owns the model and the cryptography together---and the V\&V framework here is how that stack is qualified to run sovereign and disconnected, where the data lives, rather than rented from a black box.

% ============================================================================
\section{Observability Platform}
% ============================================================================

\subsection{Architecture}

Hanzo O11y provides unified observability across three telemetry signals:

\begin{table}[H]
\centering
\begin{tabular}{@{}lp{5cm}l@{}}
\toprule
\textbf{Signal} & \textbf{Capabilities} & \textbf{Storage} \\
\midrule
Logs & Full-text search, structured filters, correlation & Datastore \\
Metrics & Time-series, histograms, percentiles, Apdex & Datastore \\
Traces & Flamegraphs, Gantt charts, span detail & Datastore \\
\bottomrule
\end{tabular}
\caption{O11y three-signal observability architecture.}
\label{tab:o11y-signals}
\end{table}

\subsection{LLM-Specific Observability}

AI workloads require specialized telemetry beyond traditional APM:

\begin{lstlisting}[caption={LLM span attributes in O11y trace.}]
{
  "name": "llm.completion",
  "duration_nano": 1500000000,
  "attrs": {
    "llm.model": "zen4-pro-27b",
    "llm.token_input": "2048",
    "llm.token_output": "512",
    "llm.temperature": "0.7",
    "llm.cost": "0.0035",
    "llm.provider": "hanzo-engine"
  }
}
\end{lstlisting}

\begin{itemize}
\item \textbf{Token tracking}: Input/output token counts per request, aggregated by model, provider, and team.
\item \textbf{Cost analysis}: Real-time cost tracking with budget alerts per organization.
\item \textbf{Quality metrics}: BLEU, ROUGE, exact match scores via evaluation workflows.
\item \textbf{Latency breakdown}: Time-to-first-token, tokens-per-second, queue wait time.
\item \textbf{Error analysis}: Provider failure rates, timeout patterns, content filter triggers.
\end{itemize}

\subsection{Alert and Anomaly Detection}

\begin{itemize}
\item \textbf{Threshold-based}: Static rules on any metric (latency $>$ 2s, error rate $>$ 5\%).
\item \textbf{Anomaly detection}: Statistical deviation from baseline patterns.
\item \textbf{Composite rules}: AND/OR combinations with flapper suppression.
\item \textbf{Multi-channel}: Email, Slack, PagerDuty, MS Teams, webhooks.
\item \textbf{Escalation}: Time-based severity escalation with on-call integration.
\end{itemize}

% ============================================================================
\section{Behavioral Analytics}
% ============================================================================

\subsection{Hanzo Insights Platform}

Insights provides behavioral analytics with IAM-integrated multi-tenancy:

\begin{itemize}
\item \textbf{Event ingestion}: Rust capture service for high-throughput event collection, Kafka-compatible streaming via Hanzo Stream.
\item \textbf{OLAP engine}: Datastore columnar storage (Hanzo's Apache-2.0 ClickHouse fork~\cite{clickhouse}) for sub-second queries across billions of events.
\item \textbf{Cohort segmentation}: Group users by behavior patterns, demographics, or temporal characteristics.
\item \textbf{Funnel analysis}: Multi-step conversion tracking with statistical significance testing.
\item \textbf{Retention modeling}: Cohort-based retention curves with churn prediction.
\item \textbf{Feature flags}: A/B testing with gradual rollout, kill switches, and variant-level metrics.
\item \textbf{Dashboards}: Time-series, heatmaps, pie charts, bar charts with custom SQL queries.
\end{itemize}

\subsection{LLM Analytics Module}

\begin{itemize}
\item \textbf{Provider adapters}: Unified client for 100+ LLM providers (OpenAI, Anthropic, local Hanzo Engine).
\item \textbf{Evaluation framework}: LLM-as-judge with configurable metrics and datasets.
\item \textbf{Trace clustering}: Semantic grouping of conversations using embeddings.
\item \textbf{Sentiment classification}: Automated sentiment analysis on model outputs.
\item \textbf{Model registry}: Version tracking, provider key management, deployment tagging.
\end{itemize}

% ============================================================================
\section{ML Pipeline Orchestration}
% ============================================================================

\subsection{Temporal Workflows}

Temporal provides durable execution for long-running ML operations:

\begin{table}[H]
\centering
\small
\begin{tabular}{@{}lp{5.5cm}@{}}
\toprule
\textbf{Workflow} & \textbf{Purpose} \\
\midrule
RunEvaluationWorkflow & LLM model evaluation with judge activities \\
BatchTraceSummarizationWorkflow & Summarize conversation traces in batches \\
TraceClusteringWorkflow & Cluster traces by semantic similarity \\
DailyTraceClusteringWorkflow & Scheduled daily clustering \\
ClassifySentimentWorkflow & Sentiment analysis on LLM outputs \\
\bottomrule
\end{tabular}
\caption{Temporal ML workflow catalog.}
\label{tab:temporal-workflows}
\end{table}

\begin{algorithm}[H]
\caption{LLM Evaluation Workflow}
\begin{algorithmic}[1]
\State Fetch evaluation configuration from database
\For{each trial $i = 1, \ldots, N$}
  \State Execute LLM judge activity (retry: 3 attempts, exponential backoff)
  \State Emit evaluation result event to analytics stream
  \State Update evaluation state in key-value store
\EndFor
\State Compute aggregate metrics (mean, std, percentiles)
\State Emit completion telemetry (success/error, duration)
\State \Return evaluation summary with per-trial results
\end{algorithmic}
\end{algorithm}

\subsection{Dagster DAG Pipelines}

Dagster manages ETL and data preparation workflows:

\begin{itemize}
\item \textbf{PostgreSQL $\to$ Datastore ETL}: Real-time event ingestion with sensor triggers.
\item \textbf{Historical backfill}: Time-partitioned asset materialization for training data.
\item \textbf{Materialized columns}: Pre-compute expensive aggregations for query performance.
\item \textbf{Person resolution}: User identity unification across events.
\item \textbf{Operational alerts}: Slack notifications on DAG failures.
\end{itemize}

\subsection{Experiment Tracking}

\begin{itemize}
\item Model comparison with A/B testing framework.
\item Statistical significance (t-test, chi-square) for metric differences.
\item Prompt version tracking with performance diffs.
\item Feature flags for gradual model rollout.
\end{itemize}

% ============================================================================
\section{Digital Twin Environments}
% ============================================================================

\subsection{Netrunner: Network Simulation}

Netrunner provides configurable network simulation for system qualification:

\begin{itemize}
\item Configurable validator count, network topology, and fault injection.
\item Consensus protocol selection (Quasar, Wave, Photon, Nova, Nebula).
\item Byzantine validator simulation for adversarial testing.
\item Performance benchmarking under controlled conditions.
\end{itemize}

\subsection{Kubernetes-Based Environment Cloning}

\begin{itemize}
\item \textbf{Namespace isolation}: Each test environment in a dedicated K8s namespace.
\item \textbf{Parallel execution}: Multiple test environments run simultaneously.
\item \textbf{Docker Compose stacks}: Dev stack (single node) and Universe (5-validator production replica).
\item \textbf{Multi-cluster}: hanzo-k8s (DigitalOcean) and apps-gke (GCP ARM64).
\end{itemize}

\subsection{Virtual Machine Registry}

38 registered VMs with pluggable lifecycle:

\begin{table}[H]
\centering
\small
\begin{tabular}{@{}llp{4cm}@{}}
\toprule
\textbf{VM} & \textbf{Chain} & \textbf{Test Focus} \\
\midrule
PlatformVM & P-Chain & Staking, validator set management \\
ExchangeVM & X-Chain & UTXO asset trading \\
EVM & C-Chain & Smart contracts, precompiles \\
DexVM & D-Chain & Orderbook, perpetuals, AMM \\
ThresholdVM & T-Chain & FHE, threshold signing \\
SessionVM & S-Chain & PQ messaging \\
BridgeVM & B-Chain & Cross-chain verification \\
AIvm & A-Chain & AI agent consensus \\
QuantumVM & Q-Chain & PQ key operations \\
\bottomrule
\end{tabular}
\caption{Selected VMs from 38 registered (security-relevant subset).}
\label{tab:vms}
\end{table}

% ============================================================================
\section{Automated Testing Frameworks}
% ============================================================================

\subsection{Test Coverage}

\begin{table}[H]
\centering
\begin{tabular}{@{}lrl@{}}
\toprule
\textbf{Component} & \textbf{Tests} & \textbf{Status} \\
\midrule
Quasar consensus & 28+ & All passing \\
Corona threshold & 21 & All passing \\
SessionVM E2E & 6 & All passing \\
DEX precompiles & 68 & All passing \\
NIST PQ vectors & Per-algorithm & All passing \\
MPC keygen/signing & 33 files & All passing \\
\bottomrule
\end{tabular}
\caption{Automated test suite coverage.}
\label{tab:test-coverage}
\end{table}

\subsection{CI/CD Pipeline}

\begin{itemize}
\item \textbf{Multi-platform builds}: \texttt{linux/amd64} + \texttt{linux/arm64} via native runners (no QEMU).
\item \textbf{ARC scale sets}: Up to 30 ARM64 runners on GKE T2A Ampere nodes.
\item \textbf{Tags}: \texttt{:main}, \texttt{:test}, \texttt{:dev} (branch-based), \texttt{:latest} (default branch).
\item \textbf{Shared workflow}: \texttt{hanzoai/.github/.github/workflows/docker-build.yml@main}.
\end{itemize}

% ============================================================================
\section{Formal Verification}
% ============================================================================

\subsection{Cryptographic Proofs}

17 formal proofs cover cryptographic primitives, consensus protocols, and cross-chain messaging:

\begin{itemize}
\item \textbf{proof-crypto-bls}: BLS12-381 aggregate signature unforgeability.
\item \textbf{proof-crypto-cggmp21}: CGGMP21 threshold ECDSA correctness under UC framework.
\item \textbf{proof-crypto-mldsa}: ML-DSA EU-CMA reduction to Module-LWE hardness.
\item \textbf{proof-crypto-tfhe}: TFHE semantic security under LWE assumption.
\item \textbf{proof-crypto-verkle}: Verkle tree binding and hiding properties.
\item \textbf{proof-protocol-wave}: Wave consensus liveness and safety.
\item \textbf{proof-protocol-nova}: Nova linear consensus termination bound.
\item \textbf{proof-bridge-teleport}: Cross-chain message integrity guarantee.
\item \textbf{proof-warp-delivery}: Warp messaging delivery with finality binding.
\end{itemize}

\subsection{Gas Determinism}

All EVM precompiles provide O(1) gas cost computation from input header bytes:

\begin{theorem}[Gas Determinism]
For any precompile $P$ with input $x$, the gas cost $G(P, x)$ is computable in constant time from the first $k$ bytes of $x$ (where $k \leq 8$), independent of the remaining input. No unbounded computation occurs during gas estimation.
\end{theorem}

% ============================================================================
\section{Performance Measurement}
% ============================================================================

\subsection{Consensus Performance}

\begin{table}[H]
\centering
\begin{tabular}{@{}lrrrr@{}}
\toprule
\textbf{Network} & \textbf{Validators} & \textbf{Finality} & \textbf{TPS} & \textbf{Config} \\
\midrule
Mainnet & 21 & 9.63 s & 1,091 & $K$=21, $\alpha$=15 \\
Testnet & 11 & 6.3 s & 1,094 & $K$=11, $\alpha$=8 \\
Local & 5 & 3.69 s & 438--1,094 & $K$=5, $\alpha$=3 \\
\bottomrule
\end{tabular}
\caption{Consensus performance across deployment sizes.}
\label{tab:consensus-perf}
\end{table}

\subsection{Wire Protocol}

\begin{table}[H]
\centering
\begin{tabular}{@{}lrrrr@{}}
\toprule
\textbf{Operation} & \textbf{ZAP} & \textbf{MCP JSON-RPC} & \textbf{Speedup} & \textbf{Allocs} \\
\midrule
Single tool call & 199 ns & 3,939 ns & 20$\times$ & 2 vs.\ 58 \\
20-tool orchestration & 5,104 ns & 73,685 ns & 14$\times$ & 60 vs.\ 1,060 \\
Parse (buffer read) & 2.2 ns & --- & --- & 0 \\
Build (message construct) & 38.9 ns & --- & --- & 0 \\
Consensus round & 4.6 ns & --- & --- & 0 \\
\bottomrule
\end{tabular}
\caption{ZAP benchmarks (Apple M1 Max, Go 1.26.1). MCP comparison from measured test suite.}
\label{tab:wire-perf}
\end{table}

% ============================================================================
\section{Data Flow Architecture}
% ============================================================================

\begin{figure}[H]
\centering
\begin{tikzpicture}[
  node distance=1.2cm,
  box/.style={draw, rectangle, minimum width=3cm, minimum height=0.8cm, align=center, font=\small},
  arrow/.style={-{Stealth[length=3mm]}, thick}
]
\node[box] (app) {Application Layer\\(LLM calls, services)};
\node[box, below=of app] (otel) {OpenTelemetry SDK\\(traces, metrics, logs)};
\node[box, below=of otel] (collect) {OTel Collector\\$\to$ Kafka $\to$ Datastore};
\node[box, below left=1cm and -0.5cm of collect] (o11y) {O11y\\Query Service};
\node[box, below=of collect] (temporal) {Temporal\\ML Workflows};
\node[box, below right=1cm and -0.5cm of collect] (insights) {Insights\\Analytics};
\node[box, right=2cm of temporal] (dagster) {Dagster\\Pipelines};

\draw[arrow] (app) -- (otel);
\draw[arrow] (otel) -- (collect);
\draw[arrow] (collect) -- (o11y);
\draw[arrow] (collect) -- (temporal);
\draw[arrow] (collect) -- (insights);
\draw[arrow] (collect) -- (dagster);
\end{tikzpicture}
\caption{V\&V data flow from application to analytics.}
\label{fig:dataflow}
\end{figure}

% ============================================================================
\section{Naval Application Scenarios}
% ============================================================================

\subsection{System Qualification}

Digital twin environments replicate production configurations for pre-deployment testing:
\begin{itemize}
\item Deploy identical K8s stack in isolated namespace.
\item Run full consensus with simulated network conditions (latency, jitter, packet loss).
\item Execute automated test suites against all cryptographic primitives.
\item Verify performance meets operational requirements before fleet deployment.
\end{itemize}

\subsection{Mission Rehearsal}

\begin{itemize}
\item Configure Netrunner with theater-specific network topology.
\item Simulate DDIL conditions (satellite jamming, link degradation).
\item Validate AI model performance under constrained bandwidth.
\item Test graceful degradation and recovery procedures.
\end{itemize}

\subsection{Continuous Monitoring at Sea}

\begin{itemize}
\item O11y agents deployed on each node collect telemetry locally.
\item Store-and-forward synchronization when connectivity resumes.
\item Anomaly detection runs locally with pre-trained models.
\item Alert escalation via available communication channels.
\end{itemize}

% ============================================================================
\section{Conclusion}
% ============================================================================

We have presented an automated V\&V framework that provides continuous verification of AI/ML systems, cryptographic implementations, and consensus protocols through integrated observability, behavioral analytics, ML pipeline orchestration, and digital twin environments. The framework's 17 formal proofs, comprehensive test suites, and production observability stack enable data-driven decision-making for naval system qualification and mission validation.

% ============================================================================
\begin{thebibliography}{18}

\bibitem{otel} OpenTelemetry, ``Observability Framework for Cloud-Native Software,'' CNCF, 2024.

\bibitem{clickhouse} ClickHouse Inc., ``ClickHouse: Fast Open-Source OLAP Database Management System,'' 2024.

\bibitem{temporal} Temporal Technologies, ``Temporal: Durable Execution Platform,'' 2024.

\bibitem{dagster} Elementl, ``Dagster: Cloud-Native Data Orchestration,'' 2024.

\bibitem{o11y} Hanzo Team, ``Hanzo O11y: Unified Observability Platform,'' Hanzo Industries, 2024.

\bibitem{insights} Hanzo Team, ``Hanzo Insights: Behavioral Analytics Platform,'' Hanzo Industries, 2024.

\bibitem{prometheus} Prometheus Authors, ``Prometheus: Monitoring System and Time Series Database,'' CNCF, 2024.

\bibitem{k8s} Kubernetes Authors, ``Kubernetes: Production-Grade Container Orchestration,'' CNCF, 2024.

\bibitem{quasar} Hanzo Team, ``Quasar Consensus,'' Hanzo, 2025.

\bibitem{corona} ``Practical Two-Round Threshold Signature from LWE,'' IACR ePrint 2024/1113.

\bibitem{fips204} NIST, ``Module-Lattice-Based Digital Signature Standard,'' FIPS 204, 2024.

\bibitem{nistpqc} NIST, ``Post-Quantum Cryptography Standardization,'' 2024.

\bibitem{digitaltwin} M. Grieves, ``Digital Twin: Manufacturing Excellence through Virtual Factory Replication,'' 2015.

\bibitem{devops} G. Kim \etal, ``The DevOps Handbook,'' IT Revolution Press, 2016.

\bibitem{abtesting} R. Kohavi \etal, ``Trustworthy Online Controlled Experiments,'' Cambridge University Press, 2020.

\end{thebibliography}

\section*{Reproducibility}
\textbf{Source code.} \texttt{github.com/hanzoai/platform, hanzoai/agents} (commit \texttt{TBD}).\\
\textbf{Build.} make test per component; CI via GitHub Actions on amd64/arm64 runners.\\
\textbf{Hardware tested.} Linux amd64/arm64; macOS for local dev.\\
\textbf{Standards referenced.} IEEE 29119, NIST SP 800-160, DevSecOps.\\
\textbf{Contact.} research@hanzo.ai.

\end{document}
